% Figure (fig:weight). Body: _generated/weight-fig.tex, written by
% pipeline/figures/plot_edge_weight.py. The fixed-cell counts drawn are
% asserted against gcd(rows-1, cols-1)+1 inside the generator, so a wrong
% picture fails `make defs` rather than printing.
\begin{figure}[!t]
\centering
\begin{tikzpicture}
\node[inner sep=2pt] (figweight) {% GENERATED by pipeline/figures/plot_edge_weight.py -- do not edit by hand.
% Fixed-cell counts recomputed from gcd(rows-1, cols-1)+1; matches src/_plan.py.
\begin{tikzpicture}[x=1pt,y=1pt]
\definecolor{dink}{HTML}{16202A}
\definecolor{dframe}{HTML}{3E4C59}
\definecolor{dground}{HTML}{F4F8FB}
\definecolor{dpre}{HTML}{DCE9F2}
\definecolor{dleft}{HTML}{A9CFE3}
\definecolor{dright}{HTML}{6FA8C9}
\definecolor{dpost}{HTML}{CFE0EC}
\definecolor{dband}{HTML}{2E5A75}
\definecolor{daccent}{HTML}{C2571A}
\definecolor{daccentpale}{HTML}{F6E2D2}
\useasboundingbox (0,0) rectangle (241.0,120.0);
\draw[draw=dframe,line width=0.7pt,rounded corners=2.4pt] (0,0) rectangle (241.0,120.0);
\draw[fill=dground,draw=dframe,line width=0.5pt,rounded corners=1.6pt] (6.00,6.00) rectangle (118.50,114.00);
\node[anchor=north west,text=dink,align=center,font=\fontsize{6}{7}\selectfont\bfseries] at (10.00,110.00) {(a)};
\draw[fill=daccentpale,draw=dink,line width=0.35pt] (42.25,76.00) rectangle (62.25,96.00);
\node[anchor=center,text=daccent,align=center,font=\fontsize{4.6}{5.4}\selectfont] at (52.25,86.00) {\textbf{0}};
\draw[fill=white,draw=dink,line width=0.35pt] (62.25,76.00) rectangle (82.25,96.00);
\node[anchor=center,text=dink,align=center,font=\fontsize{4.6}{5.4}\selectfont] at (72.25,86.00) {1};
\draw[fill=white,draw=dink,line width=0.35pt] (42.25,56.00) rectangle (62.25,76.00);
\node[anchor=center,text=dink,align=center,font=\fontsize{4.6}{5.4}\selectfont] at (52.25,66.00) {2};
\draw[fill=daccentpale,draw=dink,line width=0.35pt] (62.25,56.00) rectangle (82.25,76.00);
\node[anchor=center,text=daccent,align=center,font=\fontsize{4.6}{5.4}\selectfont] at (72.25,66.00) {\textbf{3}};
\node[anchor=center,text=dink,align=center,font=\fontsize{6}{7}\selectfont] at (62.25,105.00) {a $2\!\times\!2$ block};
\draw[fill=dpre,draw=dink,line width=0.5pt,rounded corners=1.2pt] (11.00,11.00) rectangle (113.50,35.00);
\node[anchor=center,text=dink,align=center,font=\fontsize{4.6}{5.4}\selectfont] at (62.25,28.50) {$\gcd(2\!-\!1,2\!-\!1)+1 = 2$ stay put};
\node[anchor=center,text=dink,align=center,font=\fontsize{5.2}{6.2}\selectfont] at (62.25,18.00) {$moved = 1-\tfrac{2}{4} = \mathbf{0.50}$};
\draw[fill=dground,draw=dframe,line width=0.5pt,rounded corners=1.6pt] (122.50,6.00) rectangle (235.00,114.00);
\node[anchor=north west,text=dink,align=center,font=\fontsize{6}{7}\selectfont\bfseries] at (126.50,110.00) {(b)};
\draw[fill=daccentpale,draw=dink,line width=0.35pt] (150.75,82.00) rectangle (164.75,96.00);
\node[anchor=center,text=daccent,align=center,font=\fontsize{4.6}{5.4}\selectfont] at (157.75,89.00) {\textbf{0}};
\draw[fill=white,draw=dink,line width=0.35pt] (164.75,82.00) rectangle (178.75,96.00);
\node[anchor=center,text=dink,align=center,font=\fontsize{4.6}{5.4}\selectfont] at (171.75,89.00) {1};
\draw[fill=white,draw=dink,line width=0.35pt] (178.75,82.00) rectangle (192.75,96.00);
\node[anchor=center,text=dink,align=center,font=\fontsize{4.6}{5.4}\selectfont] at (185.75,89.00) {2};
\draw[fill=white,draw=dink,line width=0.35pt] (192.75,82.00) rectangle (206.75,96.00);
\node[anchor=center,text=dink,align=center,font=\fontsize{4.6}{5.4}\selectfont] at (199.75,89.00) {3};
\draw[fill=white,draw=dink,line width=0.35pt] (150.75,68.00) rectangle (164.75,82.00);
\node[anchor=center,text=dink,align=center,font=\fontsize{4.6}{5.4}\selectfont] at (157.75,75.00) {4};
\draw[fill=daccentpale,draw=dink,line width=0.35pt] (164.75,68.00) rectangle (178.75,82.00);
\node[anchor=center,text=daccent,align=center,font=\fontsize{4.6}{5.4}\selectfont] at (171.75,75.00) {\textbf{5}};
\draw[fill=white,draw=dink,line width=0.35pt] (178.75,68.00) rectangle (192.75,82.00);
\node[anchor=center,text=dink,align=center,font=\fontsize{4.6}{5.4}\selectfont] at (185.75,75.00) {6};
\draw[fill=white,draw=dink,line width=0.35pt] (192.75,68.00) rectangle (206.75,82.00);
\node[anchor=center,text=dink,align=center,font=\fontsize{4.6}{5.4}\selectfont] at (199.75,75.00) {7};
\draw[fill=white,draw=dink,line width=0.35pt] (150.75,54.00) rectangle (164.75,68.00);
\node[anchor=center,text=dink,align=center,font=\fontsize{4.6}{5.4}\selectfont] at (157.75,61.00) {8};
\draw[fill=white,draw=dink,line width=0.35pt] (164.75,54.00) rectangle (178.75,68.00);
\node[anchor=center,text=dink,align=center,font=\fontsize{4.6}{5.4}\selectfont] at (171.75,61.00) {9};
\draw[fill=daccentpale,draw=dink,line width=0.35pt] (178.75,54.00) rectangle (192.75,68.00);
\node[anchor=center,text=daccent,align=center,font=\fontsize{4.6}{5.4}\selectfont] at (185.75,61.00) {\textbf{10}};
\draw[fill=white,draw=dink,line width=0.35pt] (192.75,54.00) rectangle (206.75,68.00);
\node[anchor=center,text=dink,align=center,font=\fontsize{4.6}{5.4}\selectfont] at (199.75,61.00) {11};
\draw[fill=white,draw=dink,line width=0.35pt] (150.75,40.00) rectangle (164.75,54.00);
\node[anchor=center,text=dink,align=center,font=\fontsize{4.6}{5.4}\selectfont] at (157.75,47.00) {12};
\draw[fill=white,draw=dink,line width=0.35pt] (164.75,40.00) rectangle (178.75,54.00);
\node[anchor=center,text=dink,align=center,font=\fontsize{4.6}{5.4}\selectfont] at (171.75,47.00) {13};
\draw[fill=white,draw=dink,line width=0.35pt] (178.75,40.00) rectangle (192.75,54.00);
\node[anchor=center,text=dink,align=center,font=\fontsize{4.6}{5.4}\selectfont] at (185.75,47.00) {14};
\draw[fill=daccentpale,draw=dink,line width=0.35pt] (192.75,40.00) rectangle (206.75,54.00);
\node[anchor=center,text=daccent,align=center,font=\fontsize{4.6}{5.4}\selectfont] at (199.75,47.00) {\textbf{15}};
\node[anchor=center,text=dink,align=center,font=\fontsize{6}{7}\selectfont] at (178.75,105.00) {a $4\!\times\!4$ block};
\draw[fill=dpre,draw=dink,line width=0.5pt,rounded corners=1.2pt] (127.50,11.00) rectangle (230.00,35.00);
\node[anchor=center,text=dink,align=center,font=\fontsize{4.6}{5.4}\selectfont] at (178.75,28.50) {$\gcd(4\!-\!1,4\!-\!1)+1 = 4$ stay put};
\node[anchor=center,text=dink,align=center,font=\fontsize{5.2}{6.2}\selectfont] at (178.75,18.00) {$moved = 1-\tfrac{4}{16} = \mathbf{0.75}$};
\end{tikzpicture}
};
\draw[black,line width=0.5pt,rounded corners=2pt]
  (figweight.south west) rectangle (figweight.north east);
\end{tikzpicture}
\caption{\capheadline{Where an edge's weight comes from.}
Shaded cells are already in the right place. A transpose does not move them, so
we do not either, and a copy has to write them anyway. The smaller the block,
the larger the share that stays put. A
$\modelMinMovedRows \times \modelMinMovedCols$ block is the best any single
swap can do.}
\label{fig:weight}
\end{figure}
